\section{AHSME 1973}
        \begin{enumerate}
        	\item (\textit{AHSME 1973}) A chord which is the perpendicular bisector of a radius of length 12 in a circle, has length 


  $\textbf{(A)}\ 3\sqrt3\qquad\textbf{(B)}\ 27\qquad\textbf{(C)}\ 6\sqrt3\qquad\textbf{(D)}\ 12\sqrt3\qquad\textbf{(E)}\ \text{ none of these}$


	\item (\textit{AHSME 1973}) One thousand unit cubes are fastened together to form a large cube with edge length 10 units; this is painted and then separated into the original cubes. The number of these unit cubes which have at least one face painted is 


  $\textbf{(A)}\ 600\qquad\textbf{(B)}\ 520\qquad\textbf{(C)}\ 488\qquad\textbf{(D)}\ 480\qquad\textbf{(E)}\ 400$


	\item (\textit{AHSME 1973}) The stronger Goldbach conjecture states that any even integer greater than 7 can be written as the sum of two different prime numbers. For such representations of the even number 126, the largest possible difference between the two primes is 


  $\textbf{(A)}\ 112\qquad\textbf{(B)}\ 100\qquad\textbf{(C)}\ 92\qquad\textbf{(D)}\ 88\qquad\textbf{(E)}\ 80$


	\item (\textit{AHSME 1973}) Two congruent 30-60-90 are placed so that they overlap partly and their hypotenuses coincide. If the hypotenuse of each triangle is 12, the area common to both triangles is 


  $\textbf{(A)}\ 6\sqrt3\qquad\textbf{(B)}\ 8\sqrt3\qquad\textbf{(C)}\ 9\sqrt3\qquad\textbf{(D)}\ 12\sqrt3\qquad\textbf{(E)}\ 24$


	\item (\textit{AHSME 1973}) Of the following five statements, I to V, about the binary operation of averaging (arithmetic mean), 


 $\text{I. Averaging is associative }$
$\text{II. Averaging is commutative }$
$\text{III. Averaging distributes over addition }$
$\text{IV. Addition distributes over averaging }$
$\text{V. Averaging has an identity element }$


 those which are always true are 


 $\textbf{(A)}\ \text{All}\qquad\textbf{(B)}\ \text{I and II only}\qquad\textbf{(C)}\ \text{II and III only}$
$\textbf{(D)}\ \text{II and IV only}\qquad\textbf{(E)}\ \text{II and V only}$


	\item (\textit{AHSME 1973}) If 554 is the base $b$ representation of the square of the number whose base $b$ representation is 24, then $b$, when written in base 10, equals 


 $\textbf{(A)}\ 6\qquad\textbf{(B)}\ 8\qquad\textbf{(C)}\ 12\qquad\textbf{(D)}\ 14\qquad\textbf{(E)}\ 16$


	\item (\textit{AHSME 1973}) The sum of all integers between 50 and 350 which end in 1 is 


  $\textbf{(A)}\ 5880\qquad\textbf{(B)}\ 5539\qquad\textbf{(C)}\ 5208\qquad\textbf{(D)}\ 4877\qquad\textbf{(E)}\ 4566$


	\item (\textit{AHSME 1973}) If 1 pint of paint is needed to paint a statue 6 ft. high, then the number of pints it will take to paint (to the same thickness) 540 statues similar to the original but only 1 ft. high is 


 $\textbf{(A)}\ 90\qquad\textbf{(B)}\ 72\qquad\textbf{(C)}\ 45\qquad\textbf{(D)}\ 30\qquad\textbf{(E)}\ 15$


	\item (\textit{AHSME 1973}) In  $\triangle ABC$ with right angle at $C$, altitude $CH$ and median $CM$ trisect the right angle. If the area of $\triangle CHM$ is $K$, then the area of $\triangle ABC$ is 


  $\textbf{(A)}\ 6K\qquad\textbf{(B)}\ 4\sqrt3\ K\qquad\textbf{(C)}\ 3\sqrt3\ K\qquad\textbf{(D)}\ 3K\qquad\textbf{(E)}\ 4K$


	\item (\textit{AHSME 1973}) If $n$ is a real number, then the simultaneous system 


 $nx+y = 1$
$ny+z = 1$
$x+nz = 1$


 has no solution if and only if $n$ is equal to 


 $\textbf{(A)}\ -1\qquad\textbf{(B)}\ 0\qquad\textbf{(C)}\ 1\qquad\textbf{(D)}\ 0\text{ or }1\qquad\textbf{(E)}\ \frac{1}2$


	\item (\textit{AHSME 1973}) A circle with a circumscribed and an inscribed square centered at the origin  of a rectangular coordinate system with positive  and  axes  and  is shown in each figure  to  below.


 
                    \begin{figure}[H]
                        \centering
                        \includegraphics[width=0.5\textwidth]{Problem/AHSME_I_1973_11_Problem_Figure_1.png}
                    \end{figure}
                    


 The inequalities


 
 \[|x|+|y|\leq\sqrt{2(x^{2}+y^{2})}\leq 2\mbox{Max}(|x|, |y|)\]


 are represented geometrically* by the figure numbered


 $\textbf{(A)}\ I\qquad\textbf{(B)}\ II\qquad\textbf{(C)}\ III\qquad\textbf{(D)}\ IV\qquad\textbf{(E)}\ \mbox{none of these}$


 * An inequality of the form $f(x, y) \leq g(x, y)$, for all $x$ and $y$ is represented geometrically by a figure showing the containment


 $\{\mbox{The set of points }(x, y)\mbox{ such that }g(x, y) \leq a\} \subset\\ \{\mbox{The set of points }(x, y)\mbox{ such that }f(x, y) \leq a\}$


 for a typical real number $a$.	


 

	\item (\textit{AHSME 1973}) The average (arithmetic mean) age of a group consisting of doctors and lawyers in 40. If the doctors average 35 and the lawyers 50 years old, then the ratio of the numbers of doctors to the number of lawyers is 


 $\textbf{(A)}\ 3: 2\qquad\textbf{(B)}\ 3: 1\qquad\textbf{(C)}\ 2: 3\qquad\textbf{(D)}\ 2: 1\qquad\textbf{(E)}\ 1: 2$


	\item (\textit{AHSME 1973}) The fraction $\frac{2(\sqrt2+\sqrt6)}{3\sqrt{2+\sqrt3}}$ is equal to 


 $\textbf{(A)}\ \frac{2\sqrt2}{3} \qquad \textbf{(B)}\ 1 \qquad \textbf{(C)}\ \frac{2\sqrt3}3 \qquad \textbf{(D)}\ \frac43 \qquad \textbf{(E)}\ \frac{16}{9}$


	\item (\textit{AHSME 1973}) Each valve $A$, $B$, and $C$, when open, releases water into a tank at its own constant rate. With all three valves open, the tank fills in 1 hour, with only valves $A$ and $C$ open it takes 1.5 hours, and with only valves $B$ and $C$ open it takes 2 hours. The number of hours required with only valves $A$ and $B$ open is 


  $\textbf{(A)}\ 1.1\qquad\textbf{(B)}\ 1.15\qquad\textbf{(C)}\ 1.2\qquad\textbf{(D)}\ 1.25\qquad\textbf{(E)}\ 1.75$


	\item (\textit{AHSME 1973}) A sector with acute central angle $\theta$ is cut from a circle of radius 6. The radius of the circle circumscribed about the sector is 


 $\textbf{(A)}\ 3\cos\theta \qquad \textbf{(B)}\ 3\sec\theta \qquad \textbf{(C)}\ 3 \cos \frac12 \theta \qquad \textbf{(D)}\ 3 \sec \frac12 \theta \qquad \textbf{(E)}\ 3$


	\item (\textit{AHSME 1973}) If the sum of all the angles except one of a convex polygon is $2190^{\circ}$, then the number of sides of the polygon must be 


 $\textbf{(A)}\ 13 \qquad \textbf{(B)}\ 15 \qquad \textbf{(C)}\ 17 \qquad \textbf{(D)}\ 19 \qquad \textbf{(E)}\ 21$


	\item (\textit{AHSME 1973}) If $\theta$ is an acute angle and $\sin \frac12 \theta = \sqrt{\frac{x-1}{2x}}$, then $\tan \theta$ equals 


 $\textbf{(A)}\ x \qquad \textbf{(B)}\ \frac1{x} \qquad \textbf{(C)}\ \frac{\sqrt{x-1}}{x+1} \qquad \textbf{(D)}\ \frac{\sqrt{x^2-1}}{x} \qquad \textbf{(E)}\ \sqrt{x^2-1}$


	\item (\textit{AHSME 1973}) If $p \geq 5$ is a prime number, then $24$ divides $p^2 - 1$ without remainder 


 $\textbf{(A)}\ \text{never} \qquad \textbf{(B)}\ \text{sometimes only} \qquad \textbf{(C)}\ \text{always} \qquad$
$\textbf{(D)}\ \text{only if } p =5 \qquad \textbf{(E)}\ \text{none of these}$


	\item (\textit{AHSME 1973}) Define $n_a!$ for $n$ and $a$ positive to be 


 
 \[n_a ! = n (n-a)(n-2a)(n-3a)...(n-ka)\]


 where $k$ is the greatest integer for which $n>ka$. Then the quotient $72_8!/18_2!$ is equal to 


 $\textbf{(A)}\ 4^5 \qquad \textbf{(B)}\ 4^6 \qquad \textbf{(C)}\ 4^8 \qquad \textbf{(D)}\ 4^9 \qquad \textbf{(E)}\ 4^{12}$


	\item (\textit{AHSME 1973}) A cowboy is 4 miles south of a stream which flows due east. He is also 8 miles west and 7 miles north of his cabin. He wishes to water his horse at the stream and return home. The shortest distance (in miles) he can travel and accomplish this is 


 $\textbf{(A)}\ 4+\sqrt{185} \qquad \textbf{(B)}\ 16 \qquad \textbf{(C)}\ 17 \qquad \textbf{(D)}\ 18 \qquad \textbf{(E)}\ \sqrt{32}+\sqrt{137}$


	\item (\textit{AHSME 1973}) The number of sets of two or more consecutive positive integers whose sum is 100 is


 $\textbf{(A)}\ 1 \qquad \textbf{(B)}\ 2 \qquad \textbf{(C)}\ 3 \qquad \textbf{(D)}\ 4 \qquad \textbf{(E)}\ 5$


	\item (\textit{AHSME 1973}) The set of all real solutions of the inequality 


 
 \[|x - 1| + |x + 2| < 5\]


 is 


 $\textbf{(A)}\ x \in ( - 3,2) \qquad \textbf{(B)}\ x \in ( - 1,2) \qquad \textbf{(C)}\ x \in ( - 2,1) \qquad$
$\textbf{(D)}\ x \in \left( - \frac32,\frac72\right) \qquad \textbf{(E)}\ \O \text{ (empty})$


 [Note: I updated the notation on this problem.]	


 

	\item (\textit{AHSME 1973}) There are two cards; one is red on both sides and the other is red on one side and blue on the other. The cards have the same probability (1/2) of being chosen, and one is chosen and placed on the table. If the upper side of the card on the table is red, then the probability that the under-side is also red is 


 $\textbf{(A)}\ \frac14 \qquad \textbf{(B)}\ \frac13 \qquad \textbf{(C)}\ \frac12 \qquad \textbf{(D)}\ \frac23 \qquad \textbf{(E)}\ \frac34$


	\item (\textit{AHSME 1973}) The check for a luncheon of 3 sandwiches, 7 cups of coffee and one piece of pie came to $\$3.15$. The check for a luncheon consisting of 4 sandwiches, 10 cups of coffee and one piece of pie came to $\$4.20$ at the same place. The cost of a luncheon consisting of one sandwich, one cup of coffee, and one piece of pie at the same place will come to 


 $\textbf{(A)}\ \$1.70 \qquad \textbf{(B)}\ \$1.65 \qquad \textbf{(C)}\ \$1.20 \qquad \textbf{(D)}\ \$1.05 \qquad \textbf{(E)}\ \$0.95$


	\item (\textit{AHSME 1973}) A circular grass plot 12 feet in diameter is cut by a straight gravel path 3 feet wide, one edge of which passes through the center of the plot. The number of square feet in the remaining grass area is 


 $\textbf{(A)}\ 36\pi-34 \qquad \textbf{(B)}\ 30\pi - 15 \qquad \textbf{(C)}\ 36\pi - 33 \qquad$
$\textbf{(D)}\ 35\pi - 9\sqrt3 \qquad \textbf{(E)}\ 30\pi - 9\sqrt3$


	\item (\textit{AHSME 1973}) The number of terms in an A.P. (Arithmetic Progression) is even. The sum of the odd and even-numbered terms are 24 and 30, respectively. If the last term exceeds the first by 10.5, the number of terms in the A.P. is 


 $\textbf{(A)}\ 20 \qquad \textbf{(B)}\ 18 \qquad \textbf{(C)}\ 12 \qquad \textbf{(D)}\ 10 \qquad \textbf{(E)}\ 8$


	\item (\textit{AHSME 1973}) Cars A and B travel the same distance. Car A travels half that distance at $u$ miles per hour and half at $v$ miles per hour. Car B travels half the time at $u$ miles per hour and half at $v$ miles per hour. The average speed of Car A is $x$ miles per hour and that of Car B is $y$ miles per hour. Then we always have 


 $\textbf{(A)}\ x \leq y\qquad \textbf{(B)}\ x \geq y \qquad \textbf{(C)}\ x=y \qquad \textbf{(D)}\ x<y\qquad \textbf{(E)}\ x>y$


	\item (\textit{AHSME 1973}) If $a$, $b$, and $c$ are in geometric progression (G.P.) with $1 < a < b < c$ and $n>1$ is an integer, then $\log_an$, $\log_bn$, $\log_cn$ form a sequence 


 $\textbf{(A)}\ \text{which is a G.P} \qquad$
$\textbf{(B)}\ \text{which is an arithmetic progression (A.P)} \qquad$
$\textbf{(C)}\ \text{in which the reciprocals of the terms form an A.P} \qquad$
$\textbf{(D)}\ \text{in which the second and third terms are the }n\text{th powers of the first and second respectively} \qquad$
$\textbf{(E)}\ \text{none of these}$


	\item (\textit{AHSME 1973}) Two boys start moving from the same point A on a circular track but in opposite directions. Their speeds are 5 ft. per second and 9 ft. per second. If they start at the same time and finish when they first meet at the point A again, then the number of times they meet, excluding the start and finish, is 


 $\textbf{(A)}\ 13 \qquad \textbf{(B)}\ 25 \qquad \textbf{(C)}\ 44 \qquad \textbf{(D)}\ \text{infinity} \qquad \textbf{(E)}\ \text{none of these}$


	\item (\textit{AHSME 1973}) Let $[t]$ denote the greatest integer $\leq t$ where $t \geq 0$ and $S = \{(x,y): (x-T)^2 + y^2 \leq T^2 \text{ where } T = t - [t]\}$. Then we have 


 $\textbf{(A)}\ \text{the point } (0,0) \text{ does not belong to } S \text{ for any } t \qquad$
$\textbf{(B)}\ 0 \leq \text{Area } S \leq \pi \text{ for all } t \qquad$
$\textbf{(C)}\ S \text{ is contained in the first quadrant for all } t \geq 5 \qquad$
$\textbf{(D)}\ \text{the center of } S \text{ for any } t \text{ is on the line } y=x \qquad$
$\textbf{(E)}\ \text{none of the other statements is true}$


	\item (\textit{AHSME 1973}) In the following equation, each of the letters represents uniquely a different digit in base ten: 


 
 \[(YE) \cdot (ME) = TTT\]


 The sum $E+M+T+Y$ equals 


 $\textbf{(A)}\ 19 \qquad \textbf{(B)}\ 20 \qquad \textbf{(C)}\ 21 \qquad \textbf{(D)}\ 22 \qquad \textbf{(E)}\ 24$


	\item (\textit{AHSME 1973}) The volume of a pyramid whose base is an equilateral triangle of side length 6 and whose other edges are each of length $\sqrt{15}$ is 


 $\textbf{(A)}\ 9 \qquad \textbf{(B)}\ 9/2 \qquad \textbf{(C)}\ 27/2 \qquad \textbf{(D)}\ \frac{9\sqrt3}{2} \qquad \textbf{(E)}\ \text{none of these}$


	\item (\textit{AHSME 1973}) When one ounce of water is added to a mixture of acid and water, the new mixture is $20\%$ acid. When one ounce of acid is added to the new mixture, the result is $33\frac13\%$ acid. The percentage of acid in the original mixture is


 $\textbf{(A)}\ 22\% \qquad \textbf{(B)}\ 24\% \qquad \textbf{(C)}\ 25\% \qquad \textbf{(D)}\ 30\% \qquad \textbf{(E)}\ 33\frac13 \%$


	\item (\textit{AHSME 1973}) A plane flew straight against a wind between two towns in 84 minutes and returned with that wind in 9 minutes less than it would take in still air. The number of minutes (2 answers) for the return trip was


 $\textbf{(A)}\ 54 \text{ or } 18 \qquad \textbf{(B)}\ 60 \text{ or } 15 \qquad \textbf{(C)}\ 63 \text{ or } 12 \qquad \textbf{(D)}\ 72 \text{ or } 36 \qquad \textbf{(E)}\ 75 \text{ or } 20$


	\item (\textit{AHSME 1973}) In the unit circle shown in the figure, chords $PQ$ and $MN$ are parallel to the unit radius $OR$ of the circle with center at $O$. Chords $MP$, $PQ$, and $NR$ are each $s$ units long and chord $MN$ is $d$ units long. 


 
                    \begin{figure}[H]
                        \centering
                        \includegraphics[width=0.5\textwidth]{Problem/AHSME_I_1973_35_Problem_Figure_1.png}
                    \end{figure}
                    


 Of the three equations


 
 \[\textbf{I.}\ d-s=1, \qquad \textbf{II.}\ ds=1, \qquad \textbf{III.}\ d^2-s^2=\sqrt{5}\]


 those which are necessarily true are


 $\textbf{(A)}\ \textbf{I}\ \text{only} \qquad\textbf{(B)}\ \textbf{II}\ \text{only} \qquad\textbf{(C)}\ \textbf{III}\ \text{only} \qquad\textbf{(D)}\ \textbf{I}\ \text{and}\ \textbf{II}\ \text{only} \qquad\textbf{(E)}\ \textbf{I, II}\ \text{and} \textbf{ III}$


\end{enumerate}